% Methods & Results 
% 数据来源：/Users/yu/code/code2601/TY/output/regression/


\documentclass[12pt, a4paper]{article}
\usepackage[UTF8]{ctex}
\usepackage{geometry}
\geometry{margin=2.5cm}
\usepackage{booktabs}
\usepackage{longtable}
\usepackage{multirow}
\usepackage{amsmath}
\usepackage{amssymb}
\usepackage{hyperref}
\usepackage{graphicx}
\usepackage{adjustbox}
\usepackage{threeparttable}
\usepackage{siunitx}
\usepackage[table]{xcolor}

\title{生成式 AI 是否推动了岗位技能需求综合化？}
\author{}
\date{}

\begin{document}
\maketitle

% ============================================================
\section{研究方法}
% ============================================================

\subsection{技能综合性的测量}

首先，我们对于岗位任务内容的综合性和岗位技能的综合性进行概念上的区分：前者指的是一个岗位实际执行的工作活动的多样性，即具体而言工人在岗位上做什么，是对岗位本身特征的描述。后者则指的是完成该岗位所需的技能组合的丰富程度。因此，本文直接测量的是岗位技能的综合性，即招聘广告中技能要求的综合程度；至于岗位任务内容是否同步发生重组，只能作为与本文结果相一致的一种可能解释。

要回答 “生成式 AI 是否推动了岗位技能需求综合化？” 这一问题，我们需要一种将非结构化的招聘广告文本转化为可量化的技能综合性指标的方法。本文构建了一套基于主题模型的测量框架，将非结构化的招聘广告文本转化为可量化的技能综合性指标。

在模型选择上，我们比较了 BERTopic、word2vec 和 LDA 等方法。考虑到数据规模（2680 万条招聘广告）和对连续概率分布的需求，最终选择了 LDA（Latent Dirichlet Allocation）。LDA 通过 Gibbs 采样能够高效处理大规模语料，且输出的主题概率分布天然适合作为技能综合性的度量基础\footnote{BERTopic 在语义表征上更优，但其聚类范式输出离散标签而非连续分布，不适合本文的因变量构建。}。

模型训练使用 268 万条文档（10\% 随机采样），迭代 1000 轮。主题数 $K=50$\footnote{具体主题见附录表~\ref{tab:all_topics}。} 通过网格搜索确定，对应主题一致性指标 $\text{CV}=0.597$\footnote{网格搜索范围为 $K \in \{20, 30, 40, 50, 60, 70, 80\}$，$K=50$ 在一致性与可解释性之间取得了较好的平衡。}。经人工审查，剔除了 6 个噪声主题（如福利待遇模板、公司地址描述、HR 通用话术等），最终保留有效主题 $K_{\text{clean}}=44$ 个。附录表~\ref{tab:all_topics} 列出了全部 50 个主题的代表性关键词及其 ALM 任务分类。

在推断阶段，我们做了两项改动。其一，出于计算效率考虑，采用解析折叠推断（folding-in）代替 Gibbs 采样对全量文档进行主题推断。其二，引入几何均值归一化（geometric mean normalization）以消除文档长度偏差：未经归一化时，较长的招聘广告天然具有更高的主题熵值，这种差异反映的是文本长度而非技能结构的差异。归一化后的主题推断公式为：
\begin{equation}
  P(\text{topic}_k \mid \text{doc}) \propto \alpha_k \cdot \left[\prod_{i=1}^{n} P(w_i \mid \text{topic}_k)\right]^{1/n}
\end{equation}
其中 $\alpha_k$ 为主题先验，$n$ 为文档词数。

基于主题分布，我们构建了 8 个互补的技能综合性指标。衡量分散程度的 5 个：Shannon 熵、有效主题数（$e^H$）、Rao 二次熵（Rao-Q）、显著主题数和尾部质量比；衡量集中程度的 3 个：HHI 指数、主导主题概率和 Gini 系数\footnote{Shannon 熵 $H = -\sum p_k \log p_k$；有效主题数 $= e^H$，可理解为技能要求相当于均匀分布在多少个领域；Rao-Q $= \sum_{j \neq k} d_{jk} p_j p_k$，$d_{jk}$ 为主题间 Jensen-Shannon 距离，弥补了 Shannon 熵不考虑主题间相似性的局限。HHI $= \sum p_k^2$；显著主题数为 $p_k > 1/K$ 的主题个数；尾部质量比为非主导主题概率质量之和与主导主题概率的比值。}。同时使用多项指标，是因为单一指标可能受分布假设影响，而不同数学构造下的指标若指向同一结论，结论更可信。

正文以 Shannon 熵为主报告指标，因为它对分布形状变化较为敏感，且在信息论和生态多样性文献中有成熟的理论基础。技能要求越集中于少数领域，熵值越低；分布越广，熵值越高。有效主题数 $e^H$ 是熵的指数变换，不改变岗位间的排序，但解释更直观：该岗位的技能要求大致相当于覆盖了多少个领域。两者回归结论一致，本文在报告效应量时同时使用，兼顾统计性质和直观性。

上述指标衡量的是技能需求的广度（岗位覆盖多少种技能领域），而非单项技能的深度。本文无法区分扩展边际（新增技能领域）和集约边际（既有领域内要求程度的变化），这一局限将在讨论部分说明。

本文借鉴"一个岗位即一组技能组合"（a job is a bundle of skills）的视角，将招聘广告中的技能结构视为岗位技能要求的外在表现。本文测量的是招聘需求侧技能要求结构的综合化，不直接等同于岗位内部实际任务内容的变化。

% ------------------------------------------------------------
\subsection{识别策略}
% ------------------------------------------------------------

为估计 AI 技术对岗位技能需求综合性的影响，我们采用连续处理变量的双重差分（DID）模型：
\begin{equation}
  Y_{o,i,t} = \beta(\text{AIExpo}_o \times \text{Post}_t) + \sum_{c} \overline{X}_{c,o} \times f(t) + \gamma_o + \delta_t + \theta_i + \varepsilon_{o,i,t}
\end{equation}
其中 $Y_{o,i,t}$ 为职业 $o$ 在行业 $i$、年份 $t$ 的技能综合性指标（先在单条招聘广告层面构造，再聚合到职业—行业—年份单元）。$\text{AIExpo}_o$ 为职业级别的 AI 暴露度；$\text{Post}_t$ 为 2022 年及之后的时期虚拟变量；$\gamma_o$、$\delta_t$、$\theta_i$ 分别为职业、年份和行业固定效应。回归采用加权最小二乘法（WLS），权重为各单元的岗位发布数量 $n_{\text{jobs}}$，标准误按 ISCO 职业码聚类。

处理变量 AI 暴露度来自 ILO Working Paper 140（Gmyrek et al., 2025），在 ISCO-08 四位码级别上构造，覆盖 303 个职业，由领域专家依据各职业任务中可被生成式 AI 替代或增强的潜力进行评估。本文将其视为职业层面的外生处理强度，刻画不同职业受生成式 AI 冲击的相对强弱。该指标有三个有利于识别的特征：它基于职业任务结构的理论评估而非从中国劳动市场数据反推，因此相对于因变量具有较好的外生性；横截面变异来自职业间既有任务构成的差异；且在样本期内为时不变值，不受冲击后内生调整的影响。在样本中，AI 暴露度取值范围为 $[0.18, 0.70]$，中位数 0.38，标准差 0.124。我们通过 SBERT 语义匹配将中国招聘广告的职位名称映射至 ISCO-08 编码，采用三层策略（原始职位名 $\to$ 清洗后名称 $\to$ 聚类回退），总匹配率达 95.7\%\footnote{三层匹配策略：首先使用 SBERT 对原始职位名称与 ISCO-08 中文描述进行语义相似度匹配；未匹配的使用清洗后名称重新匹配；最终未匹配的通过职位聚类多数投票法回退分配。}。

DID 的核心识别假设是平行趋势：在不存在 AI 冲击的反事实下，高暴露度和低暴露度职业的技能综合性变化趋势应当一致。但该假设在无条件设定下不成立。AI 暴露度与冲击前时期（2016--2021 年）的 Shannon 熵变化斜率显著正相关（$r = 0.31$, $p < 0.001$）。按暴露度三分位划分，低暴露组冲击前的熵变化斜率为 $-0.003$，中暴露组 $0.000$，高暴露组 $+0.004$。

这一前趋势并非 AI 的早期效应，而是反映了高暴露职业在既有特征上的系统性差异。ILO 暴露度较高的职业多为认知密集型白领岗位（会计、程序员、翻译等），在 2016--2021 年间本就更可能因数字化转型而出现技能需求的多样化。前趋势反映的是职业既有属性与综合化趋势之间的结构性相关。

据此，我们采用条件平行趋势设定，引入三个职业层面的基线特征——基线薪资（base\_salary）、基线学历要求（base\_edu）和基线经验要求（base\_exp），均取 2016--2018 年均值——以”基线特征 $\times$ 年份”交互项的形式进入模型，分别从市场价值、入门门槛和技能积累三个维度刻画职业间差异。表~\ref{tab:stepwise} 报告了逐步加入控制变量的分解过程：薪资和学历的边际贡献很小（$\beta$ 从 0.090 到 0.089 再到 0.086），加入经验要求后系数大幅衰减至 0.054，表明前趋势主要由经验积累维度的职业差异驱动。

\begin{table}[htbp]
\centering
\caption{逐步加入基线控制变量的分解（因变量：Shannon 熵）}
\label{tab:stepwise}
\begin{threeparttable}
\begin{tabular}{lcccc}
\toprule
 & (1) & (2) & (3) & (4) \\
\midrule
AIExpo $\times$ Post
  & $0.090^{***}$ & $0.089^{***}$ & $0.086^{***}$ & $0.054^{**}$ \\
  & (0.034)       & (0.033)       & (0.030)       & (0.022)      \\
\midrule
职业 FE                      & 291          & 291          & 291          & 291          \\
年份 FE                      & 10           & 10           & 10           & 10           \\
行业 FE                      & 19           & 19           & 19           & 19           \\
+ base\_salary $\times$ year &              & $\checkmark$ & $\checkmark$ & $\checkmark$ \\
+ base\_edu $\times$ year    &              &              & $\checkmark$ & $\checkmark$ \\
+ base\_exp $\times$ year    &              &              &              & $\checkmark$ \\
观测值                  & 20{,}521     & 20{,}521     & 20{,}521     & 20{,}521     \\
\bottomrule
\end{tabular}
\begin{tablenotes}[flushleft]
\small
\item 注：因变量为 Shannon 熵。每列在前一列基础上累加一个基线控制变量与年份的交互项。基线特征均为 2016--2018 年职业级均值。括号内为 ISCO 码聚类稳健标准误。$^{***}p<0.01$, $^{**}p<0.05$。
\end{tablenotes}
\end{threeparttable}
\end{table}

图~\ref{fig:pretrend} 展示了加入基线控制后事前趋势的变化。

\begin{figure}[htbp]
\centering
\includegraphics[width=0.85\textwidth]{../output/figures/pretrend_diagnostic.pdf}
\caption{事前趋势诊断：加入基线控制变量前后的比较（Shannon 熵，基期 = 2021）。空心灰圆为仅含三重固定效应的无控制规格，实心黑圆为加入基线特征 $\times$ 年份交互项后的条件平行趋势规格。误差线为 95\% 置信区间。仅展示冲击前时期（2016--2021），用于评估识别假设。}\label{fig:pretrend}
\end{figure}

无控制规格下，2016 年和 2017 年的系数显著偏离零线（$p = 0.016$ 和 $p = 0.041$）。加入基线控制后，事前系数向零收敛：2017 年从显著变为不显著（$p = 0.041 \to 0.079$），2018--2020 年均不显著（$p > 0.20$）。2016 年控制后仍然显著（$\beta = -0.078$, $p = 0.021$），绝对值从 $0.105$ 缩小至 $0.078$（26\%），基线控制有效但未能完全消除最远端的前趋势。因此，后文的 DID 估计应理解为条件平行趋势近似成立前提下的结果，本文将结合多条证据线共同评估。

控制变量使用冲击前的基线值而非时变值，以避免”坏控制变量”问题：如果薪资本身受到 AI 影响而成为中介变量，控制时变薪资将吸收部分真实因果效应。使用冲击前基线值仅控制既有的职业间差异，不干扰冲击后的变化路径（Rambachan \& Roth, 2023; Freyaldenhoven et al., 2019）。

原始数据由单条招聘广告构成。本文首先在招聘广告层面计算技能综合性指标，再将其按行业—职业—年份聚合，取单元格内招聘广告的均值作为回归因变量。回归的分析单元因此是行业—职业—年份单元格，而非单条招聘广告。最终面板包含 19 个行业、291 个 ISCO-08 四位码职业和 10 个年份（2016--2025）的交叉单元，共 20,521 个单元格；每个单元格至少包含 30 条招聘广告，对应约 2680 万条原始招聘广告。回归采用单元格内招聘广告数 $n_{\text{jobs}}$ 加权，使包含更多广告的单元格获得更大权重。除特别标注外，以下报告均基于条件平行趋势规格，并控制 ISCO 固定效应、年份固定效应、行业固定效应以及三个基线特征与年份的交互项。

% ============================================================
\section{实证结果}
% ============================================================

% ------------------------------------------------------------
\subsection{综合化趋势的初步事实}
% ------------------------------------------------------------

在进入因果识别之前，先考察一个基本事实：招聘技能需求的综合化是否确实在发生，不同暴露度的职业之间是否存在差异。

图~\ref{fig:trend_tercile} 报告了按 AI 暴露度三分位分组后，Shannon 熵和有效主题数的加权均值趋势（2016--2024 年）。两个指标的模式一致。整个样本期间内，高暴露组（T3）的技能综合性始终高于低暴露组（T1），差距在 2017 年之后基本持续。三组在 2016--2021 年间的变化方向大体一致，均经历了 2017--2018 年的上升和 2019--2020 年的回落，综合化趋势并非高暴露职业独有。但 2021 年之后出现分化：T3 继续上升，T1 明显下降，三组间差距在 2022--2024 年间持续扩大。

\begin{figure}[htbp]
\centering
\includegraphics[width=0.95\textwidth]{../output/figures/trend_by_tercile.pdf}
\caption{技能综合性的分组趋势（按 AI 暴露度三分位）。加权均值，权重为各单元岗位发布数量。虚线标记 2022 年（生成式 AI 扩散起点），灰色区域为冲击后时期。2025 年数据不完整，未纳入。}
\label{fig:trend_tercile}
\end{figure}

综合化不仅存在于高暴露职业，中低暴露组也有类似趋势，但 2022 年之后的分化提示 AI 暴露度可能与之相关。不过，描述性趋势无法识别 AI 的边际作用，因为 T3 在冲击前就持续高于 T1。下一节转向 DID 估计以分离暴露度的独立贡献。

% ------------------------------------------------------------
\subsection{AI 暴露度的边际效应}
% ------------------------------------------------------------

表~\ref{tab:main} 报告了 DID 估计结果。条件平行趋势规格下，5 个分散度指标系数全部为正，3 个集中度指标系数全部为负，均在 5\% 水平显著。相对于低暴露职业，高 AI 暴露职业的招聘广告在 2022 年之后技能需求分布更为分散，从少数领域集中转向更多领域覆盖。该估计依赖于条件平行趋势近似成立的假设，后文将评估其稳健性。

\begin{table}[htbp]
\centering
\caption{AI 暴露度对技能综合性的主效应}
\label{tab:main}
\begin{threeparttable}
\adjustbox{max width=\textwidth}{%
\begin{tabular}{l*{3}{c}@{\hspace{12pt}}*{3}{c}@{\hspace{12pt}}c}
\toprule
 & \multicolumn{3}{c}{无控制} & \multicolumn{3}{c}{条件平行趋势} & \\
\cmidrule(lr){2-4} \cmidrule(lr){5-7}
因变量 & $\beta$ & SE & $p$ & $\beta$ & SE & $p$ & 缩小 \\
\midrule
\multicolumn{8}{l}{\textit{Panel A：分散度指标（$\beta > 0$ 表示更分散/综合）}} \\[4pt]
Shannon 熵    & $0.090^{***}$  & 0.034 & 0.007 & $0.054^{**}$   & 0.022 & 0.016 & 40\% \\[3pt]
有效主题数    & $2.150^{***}$  & 0.825 & 0.009 & $1.191^{**}$   & 0.506 & 0.019 & 45\% \\[3pt]
Rao-Q         & $0.054^{***}$  & 0.020 & 0.006 & $0.035^{**}$   & 0.014 & 0.012 & 36\% \\[3pt]
显著主题数    & $1.033^{***}$  & 0.362 & 0.004 & $0.652^{**}$   & 0.253 & 0.010 & 37\% \\[3pt]
尾部质量比    & $0.070^{***}$  & 0.026 & 0.007 & $0.040^{**}$   & 0.016 & 0.013 & 44\% \\[3pt]
\midrule
\multicolumn{8}{l}{\textit{Panel B：集中度指标（$\beta < 0$ 表示更分散/综合）}} \\[4pt]
HHI           & $-0.097^{***}$ & 0.035 & 0.005 & $-0.063^{**}$  & 0.027 & 0.020 & 35\% \\[3pt]
主导主题概率  & $-0.086^{***}$ & 0.031 & 0.005 & $-0.057^{**}$  & 0.023 & 0.014 & 34\% \\[3pt]
Gini 系数     & $-0.043^{**}$  & 0.017 & 0.011 & $-0.022^{**}$  & 0.010 & 0.029 & 49\% \\[3pt]
\midrule
$N$                                & \multicolumn{3}{c}{20{,}521} & \multicolumn{3}{c}{20{,}521} & \\
Occupation FE                      & \multicolumn{3}{c}{Yes}      & \multicolumn{3}{c}{Yes}      & \\
Year FE                            & \multicolumn{3}{c}{Yes}      & \multicolumn{3}{c}{Yes}      & \\
Industry FE                        & \multicolumn{3}{c}{Yes}      & \multicolumn{3}{c}{Yes}      & \\
Baseline controls $\times$ year    & \multicolumn{3}{c}{No}       & \multicolumn{3}{c}{Yes}      & \\
Weights                            & \multicolumn{3}{c}{$n_{\text{jobs}}$} & \multicolumn{3}{c}{$n_{\text{jobs}}$} & \\
Cluster                            & \multicolumn{3}{c}{ISCO}     & \multicolumn{3}{c}{ISCO}     & \\
\bottomrule
\end{tabular}}% end adjustbox
\begin{tablenotes}[flushleft]
\small
\item 注：观测值为行业—职业—年份单元格，$N = 20{,}521$，对应约 2680 万条原始招聘广告。``无控制''规格包含职业、年份和行业固定效应；``条件平行趋势''规格在此基础上进一步加入基线薪资、基线学历要求和基线经验要求（均取 2016--2018 年职业级均值）与年份的交互项。``缩小''为从无控制到条件平行趋势后 $|\beta|$ 的衰减比例。回归按单元格内招聘广告数 $n_{\text{jobs}}$ 加权，标准误按 ISCO 职业码聚类。$^{***}p<0.01$, $^{**}p<0.05$。
\end{tablenotes}
\end{threeparttable}
\end{table}

从无控制到条件平行趋势规格，系数缩小约 34\%--49\%，均值约 40\%。无控制的 $\beta$ 可视为效应上界，因为它包含了高暴露职业基线特征上的既有优势；条件平行趋势下的 $\beta$ 可能更接近下界，因为基线薪资、学历和经验本身可能部分处于 AI 影响的传导链条中，控制它们有过度吸收的风险。约 40\% 的效应可归因于既有特征趋势，剩余约 60\% 不能被基线特征趋势解释，但也不能完全归因于 AI，因为条件平行趋势假设仍存在残余偏离。

以有效主题数衡量效应量。在 50 个 LDA 主题中，岗位的有效主题数（技能分布的等价均匀类别数）加权均值约 6.3 个（标准差 1.4），中位数 6.7 个（四分位距 5.3--8.3）。条件平行趋势规格下，AI 暴露度每增加一个标准差（$\Delta\text{AIExpo} = 0.124$），2022 年之后有效技能领域额外增加 $1.191 \times 0.124 = 0.15$ 个，约为均值的 2.4\%；以暴露度的四分位距变化（$\Delta\text{expo} = 0.190$）计算，额外增加 $1.191 \times 0.190 = 0.23$ 个，约 3.6\%\footnote{以 Shannon 熵报告，暴露度每增加一个标准差，熵值增加 $0.054 \times 0.124 = 0.007$ 个单位，约为熵标准差（0.071）的 0.094 倍。}。绝对效应量较小，但这是三重固定效应加基线趋势控制后的净增量，模型已吸收大量变异；且 8 个数学构造不同的指标方向完全一致。

% ------------------------------------------------------------
\subsection{事件研究}
% ------------------------------------------------------------

DID 估计给出了效应的平均幅度，但未揭示时间动态。图~\ref{fig:event_study} 报告了 Shannon 熵的逐年系数（基期 2021 年），事前趋势诊断已在识别策略部分完成（图~\ref{fig:pretrend}），本节聚焦冲击后一侧。

\begin{figure}[htbp]
\centering
\includegraphics[width=0.9\textwidth]{../output/figures/event_study_entropy_cond_pt.pdf}
\caption{事件研究：Shannon 熵的逐年系数（条件平行趋势规格，基期 = 2021）}
\label{fig:event_study}
\end{figure}

表~\ref{tab:event_study} 列出了各年份的精确系数，以便与图对照。

\begin{table}[htbp]
\centering
\caption{事件研究系数：Shannon 熵（基期 = 2021）}
\label{tab:event_study}
\begin{threeparttable}
\begin{tabular}{lcc}
\toprule
 & (1) & (2) \\
 & 无控制 & 条件平行趋势 \\
\midrule
2016 & $-0.105^{**}$  & $-0.078^{**}$  \\
     & (0.043)        & (0.034)        \\[2pt]
2017 & $-0.066^{**}$  & $-0.047$       \\
     & (0.032)        & (0.027)        \\[2pt]
2018 & $-0.046^{*}$   & $-0.027$       \\
     & (0.027)        & (0.026)        \\[2pt]
2019 & $-0.083$       & $-0.068$       \\
     & (0.059)        & (0.054)        \\[2pt]
2020 & $-0.065$       & $-0.057$       \\
     & (0.055)        & (0.053)        \\[2pt]
2021 & \multicolumn{2}{c}{(基期)}      \\[2pt]
2022 & $+0.017$       & $+0.007$       \\
     & (0.021)        & (0.019)        \\[2pt]
2023 & $+0.020$       & $-0.001$       \\
     & (0.020)        & (0.018)        \\[2pt]
2024 & $+0.117^{**}$  & $+0.074^{**}$  \\
     & (0.047)        & (0.036)        \\[2pt]
2025 & $-0.048$       & $-0.046$       \\
     & (0.051)        & (0.066)        \\
\bottomrule
\end{tabular}
\begin{tablenotes}[flushleft]
\small
\item 注：因变量为 Shannon 熵，基期为 2021 年。$N = 20{,}521$。2025 年数据不完整。括号内为 ISCO 码聚类稳健标准误。$^{**}p<0.05$, $^{*}p<0.10$。
\end{tablenotes}
\end{threeparttable}
\end{table}

2022 年和 2023 年的系数接近零且不显著（$\beta_{2022} = 0.007$, $p = 0.708$；$\beta_{2023} = -0.001$, $p = 0.953$），效应并非即时发生而是存在滞后。2024 年系数显著为正（$\beta = 0.074$, $p = 0.041$），与 AI 从发布到组织采纳需要时间的预期一致。替换处理期定义的检验支持这一判断：将 Post 分别定义为 $\geq$2022 和 $\geq$2023 时，熵系数均为 $+0.054$，2022 年对合并估计量的贡献基本为零，效应积累主要始于 2023--2024 年间\footnote{2024 年纯截面回归也提供一致证据：AI 暴露度与 Shannon 熵的截面相关系数为 $+0.118$（$p = 0.06$），而 2021 年同一规格下不显著（$p = 0.57$）。}。2025 年系数不显著（$\beta = -0.046$, $p = 0.486$），该年份数据尚不完整。

% ------------------------------------------------------------
\subsection{异质性分析}
% ------------------------------------------------------------

以下两组异质性分析为主效应提供支持性证据：如果 AI 暴露度的确与技能综合化相关，那么不同暴露水平和行业之间应当呈现出与理论预期一致的差异模式。

\subsubsection{暴露度阈值效应}

我们按 AI 暴露度的四分位数将职业分为 Q1--Q4 四组，以 Q1（最低暴露）为参照组，估计各分位组相对于 Q1 在冲击后时期的差异效应。结果呈现出显著的非线性模式（表~\ref{tab:quartile}）。

\begin{table}[htbp]
\centering
\caption{暴露度四分位的梯度效应}
\label{tab:quartile}
\begin{threeparttable}
\begin{tabular}{lccc}
\toprule
 & (1) & (2) & (3) \\
 & Q2 vs Q1 & Q3 vs Q1 & Q4 vs Q1 \\
\midrule
Shannon 熵  & $0.002$        & $0.023^{**}$   & $0.018^{*}$   \\
有效主题数  & $-0.001$       & $0.531^{**}$   & $0.397$       \\
显著主题数  & $0.035$        & $0.282^{**}$   & $0.224^{*}$   \\
HHI         & $-0.003$       & $-0.027^{**}$  & $-0.021^{*}$  \\
\bottomrule
\end{tabular}
\begin{tablenotes}[flushleft]
\small
\item 注：$N = 20{,}521$。报告各分位组相对于 Q1 的 DID 系数。所有回归均包含条件平行趋势控制。其余 4 个指标呈现相同模式，详见附录。$^{**}p<0.05$, $^{*}p<0.10$。
\end{tablenotes}
\end{threeparttable}
\end{table}

结果呈现非线性模式。Q2 在所有指标上效应接近零且不显著，中低暴露职业几乎未受影响。Q3 效应最强，8 个指标中 6 个在 5\% 水平显著。Q4 方向一致但略弱于 Q3，多数仅在 10\% 水平边际显著。这更像阈值效应而非线性梯度：暴露度须超过某个门槛才会显著改变技能结构。三分位分析也证实了这一点：仅 T3 系数显著（Shannon 熵：$\beta = 0.018$, $p = 0.058$），T2 不显著。

Q3 vs Q1 的效应量为 0.023 个熵单位（0.32 个标准差），远大于连续处理下暴露度增加一个标准差的平均效应（0.094 个标准差），因为有门槛效应的群体与无效应群体之间的差距，天然大于暴露度均匀推移一个标准差的平均效应。

逐一剔除分析排除了少数异常职业驱动结果的可能：每次剔除最高暴露三分位组中权重最大的一个 ISCO 职业码，$\beta$ 在 0.050--0.064 之间，全部在 $p < 0.03$ 水平显著。

\subsubsection{行业异质性}

我们进一步通过行业分组的三重交互项考察 AI 影响的行业差异（表~\ref{tab:industry}）。参照国家统计局《数字经济及其核心产业统计分类（2021）》和 OECD 知识密集型服务业分类标准，将 19 个行业分为四组：数字技术产业（信息传输、软件和信息技术服务业 I）、知识密集型服务业（金融 J、房地产 K、租赁和商务服务 L、科学研究 M、水利环保 N、教育 P、卫生 Q）、第二产业（采矿 B、制造 C、电热燃气 D、水利环境 E）和传统服务业（建筑 F、批发零售 G、交通运输 H、居民服务 O、文体娱乐 R、其他服务 S），以农林牧渔等少量剩余行业为参照组。

\begin{table}[htbp]
\centering
\caption{行业异质性：Shannon 熵的三重交互效应}
\label{tab:industry}
\begin{threeparttable}
\begin{tabular}{lcc}
\toprule
行业组 & 差异效应 $\Delta$ & 总效应 \\
\midrule
参照组（农林牧渔等） & ---                & $0.090^{***}$ \\
数字技术产业 (I)      & $-0.019$           & $0.072$       \\
知识密集型服务业      & $-0.048^{**}$      & $0.043$       \\
第二产业 (B/C/D/E)    & $-0.060^{***}$     & $0.030$       \\
传统服务业            & $-0.023$           & $0.067$       \\
\bottomrule
\end{tabular}
\begin{tablenotes}[flushleft]
\small
\item 注：$N = 20{,}521$。基础效应为参照组的 AIExpo $\times$ Post 系数；额外效应 $\Delta$ 为各行业组与参照组的差异。$^{***}p<0.01$，$^{**}p<0.05$。知识密集型服务业含 J/K/L/M/N/P/Q；传统服务业含 F/G/H/O/R/S。
\end{tablenotes}
\end{threeparttable}
\end{table}

总效应呈梯度分布：数字技术产业（0.072）和传统服务业（0.067）最高，知识密集型服务业居中（0.043），第二产业最低（0.030）。第二产业与参照组的差异高度显著（$\Delta = -0.060$, $p < 0.001$），制造业、采矿业中高暴露职业的综合化较弱，与 ALM 框架关于任务可替代性差异的预期一致\footnote{Autor, D.H., Levy, F. \& Murnane, R.J. (2003). ``The Skill Content of Recent Technological Change: An Empirical Exploration.'' \textit{The Quarterly Journal of Economics}, 118(4), 1279--1333.}。知识密集型服务业差异效应也显著为负（$\Delta = -0.048$, $p = 0.023$），可能因为这些行业技能结构起点已较多元，进一步综合化的空间有限。数字技术产业差异效应不显著（$\Delta = -0.019$, $p = 0.266$），总效应处于较高水平。

综合两组异质性分析，综合化并非均匀分布，而是呈现与理论预期一致的选择性模式：暴露度存在阈值效应，服务业和数字经济行业中的高暴露职业表现更明显。这些模式与 AI 推动综合化的假说方向一致，但属于描述性发现，不构成独立的因果识别。

% ------------------------------------------------------------
\subsection{补充证据：技能描述的内容变化}
% ------------------------------------------------------------

我们进一步考察技能描述的内容层面。将 44 个 LDA 主题依据 ALM 任务分类框架（Autor, Levy \& Murnane, 2003）映射为五类：非常规分析型（NRA, 13 个）、非常规交互型（NRI, 14 个）、常规认知型（RC, 8 个）、常规体力型（RM, 5 个）和通用/模板型（M, 14 个）\footnote{M 类主题包含”吃苦耐劳””责任心””福利待遇”等模板化、无实质技能内容的 JD 词汇。分类由两位研究者独立完成，Cohen's $\kappa > 0.8$。}，以各类别的主导主题份额为因变量，沿用同一 DID 框架估计。

M 类（通用/模板型）份额在高暴露职业中显著下降（$\beta = -0.133$, $p = 0.041$），模板化表述让位于具体技能描述\footnote{RM（常规体力型）份额显著上升（$\beta = +0.172$, $p = 0.004$），但 RM 仅含 5 个主题且粒度较粗，需谨慎解读。详见附录。}。这与熵上升方向一致，确认了熵的变化并非源于噪声或文本风格的随机波动，但描述具体化本身是测量层面的表征，不构成经济学意义上的机制。

技能综合化背后可能有两条经济机制。\textbf{任务合并}：招聘存在固定成本（搜寻、面试、培训），AI 加速了每项任务的完成，企业有动机将更多任务打包到单个岗位，降低固定成本分摊。\textbf{选择效应}：部分仅需单一技能的岗位被 AI 替代或整合，剩余岗位的技能结构因此更多元，观察到的综合化可能是岗位构成的变化而非既有岗位内部的技能扩展。本文数据无法区分两条路径，因为观察到的是职业—行业层面的聚合变化；区分它们需要企业层面的微观数据。

% ------------------------------------------------------------
\subsection{稳健性检验}
% ------------------------------------------------------------

以上结果依赖于一组特定的建模选择（处理期起点、Post 定义方式、职业分类）。表~\ref{tab:robust} 汇总了逐一替换关键设定后的结果。

\begin{table}[htbp]
\centering
\caption{稳健性检验汇总}
\label{tab:robust}
\begin{threeparttable}
\begin{tabular}{lccl}
\toprule
检验 & 显著/8 & Entropy $\beta$ & 备注 \\
\midrule
时变处理变量（累计 LLM 数） & 8/8 & $+0.195^{***}$ & 与职业趋势共线 $r\approx 0.99$ \\
+ 行业 $\times$ 年份线性趋势  & 8/8 & $+0.045^{**}$  & 效应跨行业稳健 \\
Spline DID（断点 2021）       & 0/8 & $\beta_2 = +0.009$ & 无离散结构突变 \\
ISCO-08 替代 HDBSCAN 聚类    & 8/8 & $+0.054^{**}$  & $\beta$ 缩小 20--30\% \\
Post $\geq$ 2024（替换处理期定义）& 8/8 & $+0.078^{***}$ & 当前 Post$\geq$2022 为保守估计 \\
随机化推断（500 次置换）      & --- & $+0.054^{**}$  & RI $p < 0.002$ \\
Rambachan \& Roth 敏感性        & --- & $+0.071$ (2024) & 击穿 $\bar{M} \approx 0.02$ \\
\bottomrule
\end{tabular}
\begin{tablenotes}[flushleft]
\small
\item 注：“显著/8”指 8 个因变量中在 $p<0.05$ 水平显著的个数；“熵系数”指 Shannon 熵对应的 $\text{treat}\_\text{x}\_\text{post}$ 系数；样条 DID 的 $\beta_2$ 为 2021 年后的额外斜率加速项。$^{***}p<0.01$, $^{**}p<0.05$。
\end{tablenotes}
\end{threeparttable}
\end{table}

时变处理变量（Epoch AI 累计语言模型数量）下 8 个指标全部显著，但与职业特定线性趋势共线性极强（$r \approx 0.99$），无法区分 AI 渐进渗透与既有趋势。加入行业 $\times$ 年份线性趋势后效应仍全部显著，排除了特定行业驱动的可能。样条 DID 的加速项 $\beta_2$ 在所有指标和安慰剂断点（2018、2019）上均不显著，效应更可能是渐进累积而非离散冲击。职业分类从 74 个 HDBSCAN 聚类切换到 291 个 ISCO-08 四位码后效应稳健，$\beta$ 缩小 20--30\%。处理期从 Post $\geq$ 2022 推迟至 Post $\geq$ 2024 后，Shannon 熵系数从 $+0.054$ 增至 $+0.078$（$p = 0.009$），当前定义为保守估计。随机化推断（500 次置换）的 $p < 0.002$，排除了效应源于偶然匹配。Rambachan \& Roth (2023) 敏感性分析显示 2024 年效应在标准假设下显著，但对残余前趋势偏离较为敏感（附录表~\ref{tab:app_honestdid}），单一 DID 估计量不宜作为决定性证据。

结论不依赖于 2022 年这一特定起点。处理期起点替换为 2023 或 2024 后结果方向一致且效应更强，显著效应集中在 2024 年，中国劳动力市场对生成式 AI 的调整存在滞后（附录表~\ref{tab:app_post_timing}、\ref{tab:app_placebo_timing}，图~\ref{fig:app_event_study_timing}）。

主效应在处理期定义、职业分类、行业趋势控制等维度上稳健，随机化推断排除了处理分配的偶然性，但敏感性分析显示对条件平行趋势假设依赖较高。本文不将结论完全建立在 DID 的统计显著性上，而是综合多条证据线评估：8 个指标方向完全一致且全部显著；随机化推断排除偶然匹配；非线性阈值模式（Q3 $>$ Q4 $>$ Q2 $\approx$ 0）与技术采纳门槛预期一致；行业异质性与 ALM 框架吻合；时间滞后模式与技术扩散周期相符；M 类主题份额下降提供了内容层面的支持。这些证据来自不同维度，虽非统计意义上的独立检验，但逻辑上相互补充，与``生成式 AI 很可能推动了招聘技能需求综合化''的判断方向一致。

% ============================================================
\appendix
\section{LDA 主题一览}
% ============================================================

表~\ref{tab:all_topics} 列出了全局 LDA 模型（$K=50$）的全部主题及其代表性关键词。被剔除的 6 个噪声主题以灰色标注。ALM 分类依据 Autor, Levy \& Murnane (2003) 任务框架.\footnote{Autor, D.H., Levy, F. \& Murnane, R.J. (2003). ``The Skill Content of Recent Technological Change: An Empirical Exploration.'' \textit{The Quarterly Journal of Economics}, 118(4), 1279--1333. DOI: 10.1162/003355303322552801。}。

\begin{longtable}{clp{8.5cm}l}
\caption{全部 50 个 LDA 主题及其 ALM 分类}\label{tab:all_topics} \\
\toprule
ID & 类别 & 代表性关键词 & 简称 \\
\midrule
\endfirsthead
\toprule
ID & 类别 & 代表性关键词 & 简称 \\
\midrule
\endhead
\bottomrule
\endfoot
0  & NRI & 店铺, 数据分析, 营销, 网站, 用户, 电商, 策划 & 电商运营 \\
1  & RM  & 维修, 运行, 记录, 维护保养, 物业, 现场, 检修 & 设备维护 \\
2  & NRI & 保险, 贷款, 旅游, 金融, 收入, 银行 & 保险/金融销售 \\
3  & NRA & 法律, 合同, 案件, 知识产权, 法务 & 法律事务 \\
4  & RM  & 车辆, 汽车, 维修, 物流, 司机, 驾驶, 配送 & 驾驶/物流 \\
5  & NRA & 医院, 医疗, 护理, 患者, 健康, 治疗, 医生 & 医疗护理 \\
6  & RM  & 小时, 每天, 男女不限, 宿舍, 厂区, 车间 & 工厂/体力岗 \\
7  & NRA & 投资, 金融, 交易, 风险, 审计, 融资, 银行 & 金融分析 \\
8  & NRI & 经销商, 终端, 门店, 营销, 达成, 品牌 & 渠道营销 \\
9  & M   & 积极主动, 餐厅, 认真踏实, 为人正直, 团队精神 & 通用素质 \\
10 & NRI & 链家, 房源, 咨询, 房产, 门店, 房地产 & 房产经纪 \\
11 & RC  & 人力资源, 人事, 办理, 人力资源管理, 考勤 & 人事行政 \\
12 & M   & 培养, 方向, 从事, 人才, 应届生, 热爱 & 应届生招聘 \\
13 & RM  & 施工, 工程, 现场, 验收, 规范, 材料, 图纸 & 施工管理 \\
14 & NRA & 制作, 设计师, 创意, 美术, 平面设计, 作品 & 创意设计 \\
\rowcolor[gray]{0.9}
15 & --- & 业绩, 收入, 薪酬福利, 阳光透明 & \textit{噪声：薪酬话术} \\
16 & NRI & 建设, 体系, 规划, 经营, 制度, 考核 & 体系建设/管理 \\
17 & M   & 中国, 品牌, 科技, 一家, 实现, 创新, 人才 & 公司简介 \\
18 & RC  & 订单, 记录, 客服, 客户服务, 统计, 合同, 对接 & 客服/订单 \\
19 & NRI & 测试, 客户, 产品销售, 维护客户关系, 活动策划 & 销售支持 \\
20 & NRA & 应用, 实现, 框架, 编写, 研发, 数据库, 编程 & 软件开发 \\
\rowcolor[gray]{0.9}
21 & --- & 收入, 不是, 每天, 业绩, 一起, 选择 & \textit{噪声：激励话术} \\
22 & RC  & 银行, 办理, 电话, 用户, 客服, 催收 & 银行客服/催收 \\
23 & RC  & 检查, 记录, 安全生产, 保证, 落实, 制度 & 安全管理 \\
24 & NRA & 工程, 合同, 成本, 投标, 招标, 预算 & 工程造价 \\
\rowcolor[gray]{0.9}
25 & --- & 节日, 法定, 旅游, 五险, 补助, 周末双休 & \textit{噪声：福利模板} \\
26 & NRI & 策划, 撰写, 品牌, 视频, 媒体, 文案 & 内容/新媒体 \\
27 & NRI & 学员, 课程, 校区, 咨询, 学生, 教育, 家长 & 教育咨询 \\
\rowcolor[gray]{0.9}
28 & --- & 当天安排住宿, 食宿宿舍, 叉车安全, 规章制度 & \textit{噪声：食宿/规章} \\
29 & RC  & 文件, 行政, 会议, 接待, 办公, 登记 & 行政文秘 \\
30 & NRA & 财务, 会计, 报表, 税务, 核算, 财务管理 & 财务会计 \\
31 & NRA & 测试, 调试, 应用, 工艺, 研发, 自动化, 电气 & 研发测试 \\
32 & NRA & support, responsible, design, process, business & 外企/英文岗 \\
33 & RC  & 采购, 供应商, 工艺, 物料, 品质, 成本 & 采购/供应链 \\
34 & NRI & 直播, 主播, 小时, 每天, 粉丝, 直播间 & 直播运营 \\
35 & NRA & 游戏, 基础, IT, 计算机, 热爱, 逻辑思维 & 游戏/IT \\
36 & NRI & 招商, 商户, 商家, 品牌, 拓展, 商业, 商场 & 招商拓展 \\
37 & RC  & 门店, 仓库, 商品, 顾客, 物流, 店长 & 零售/仓储 \\
38 & NRI & 营销, 市场营销, 市场信息, 销售经验, 开拓 & 市场营销 \\
\rowcolor[gray]{0.9}
39 & --- & 旅游, 提供住宿, 业绩, 费用, 封顶 & \textit{噪声：福利/业绩} \\
40 & NRA & 网络, 运维, IT, 应用, 计算机, 监控, 配置 & 网络运维 \\
41 & NRI & 电话, 思维, 关系, 开发新客户, 挖掘客户 & 电话销售 \\
42 & RM  & 人品端正, 做事仔细, 住宿, 送货, 司机 & 配送/体力岗 \\
43 & NRI & 教学, 学生, 课程, 教育, 学校, 教师, 授课 & 教育教学 \\
44 & NRA & 研发, 检测, 实验室, 研究, 实验, 报告 & 实验/研发 \\
45 & NRA & 海外, 英语, 外贸, 订单, 翻译, 国际 & 外贸/翻译 \\
\rowcolor[gray]{0.9}
46 & --- & 中心, 电话, 乘车路线, 联系电话, 街道 & \textit{噪声：地址描述} \\
47 & NRA & 解决方案, 项目管理, 推动, 规划, 数据分析 & 项目/产品管理 \\
48 & RM  & 小时, 每天, 酒店, 男女不限, 客人, 餐厅 & 酒店/餐饮 \\
\rowcolor[gray]{0.9}
49 & --- & 高铁动车, 铁路公寓, 审核合格, 客运员 & \textit{噪声：特定模板} \\
\end{longtable}

\noindent 注：类别标记——NRA: 非常规分析型, NRI: 非常规交互型, RC: 常规认知型, RM: 常规体力型, M: 通用/模板型。灰色行为剔除的噪声主题（$K_{\text{clean}} = 44$）。关键词取自各主题概率最高的前 6--8 个词。

% ============================================================
\section{处理期定义的稳健性检验}
% ============================================================

\begin{longtable}{lccc}
\caption{处理期定义的稳健性检验}\label{tab:app_post_timing} \\
\toprule
& Post $\geq$ 2022 & Post $\geq$ 2023 & Post $\geq$ 2024 \\
\midrule
\endfirsthead
\toprule
& Post $\geq$ 2022 & Post $\geq$ 2023 & Post $\geq$ 2024 \\
\midrule
\endhead
\bottomrule
\endfoot
Shannon 熵
  & $0.054^{**}$   & $0.054^{**}$   & $0.078^{***}$  \\
  & (0.022)        & (0.022)        & (0.030)        \\[4pt]
有效主题数
  & $1.191^{**}$   & $0.998^{**}$   & $1.240^{**}$   \\
  & (0.506)        & (0.446)        & (0.575)        \\[4pt]
Rao-Q
  & $0.035^{**}$   & $0.035^{**}$   & $0.052^{***}$  \\
  & (0.014)        & (0.014)        & (0.019)        \\[4pt]
显著主题数
  & $0.652^{**}$   & $0.622^{***}$  & $0.835^{***}$  \\
  & (0.253)        & (0.232)        & (0.316)        \\[4pt]
尾部质量比
  & $0.040^{**}$   & $0.034^{**}$   & $0.044^{**}$   \\
  & (0.016)        & (0.014)        & (0.019)        \\[4pt]
\midrule
\multicolumn{4}{l}{\textit{集中度指标（$\beta < 0$ 表示更分散/综合）}} \\[2pt]
HHI
  & $-0.063^{**}$  & $-0.070^{***}$ & $-0.108^{***}$ \\
  & (0.027)        & (0.027)        & (0.039)        \\[4pt]
主导主题概率
  & $-0.057^{**}$  & $-0.058^{***}$ & $-0.086^{***}$ \\
  & (0.023)        & (0.023)        & (0.032)        \\[4pt]
Gini 系数
  & $-0.022^{**}$  & $-0.018^{**}$  & $-0.023^{**}$  \\
  & (0.010)        & (0.009)        & (0.011)        \\[4pt]
\midrule
Observations       & 20{,}521 & 20{,}521 & 20{,}521 \\
$N_{\text{post}}$  & 8{,}183  & 5{,}434  & 3{,}129  \\
Occupation FE      & Yes    & Yes    & Yes    \\
Year FE            & Yes    & Yes    & Yes    \\
Industry FE        & Yes    & Yes    & Yes    \\
条件平行趋势控制 & Yes & Yes  & Yes    \\
Weights            & $n_{\text{jobs}}$ & $n_{\text{jobs}}$ & $n_{\text{jobs}}$ \\
SE clustered by ISCO & Yes  & Yes    & Yes    \\
\end{longtable}

\noindent 注：本表报告替换处理期起点后的 DID 估计结果。基准设定将 $\text{Post}_t$ 定义为 2022 年及之后，替代设定分别将其定义为 2023 年及之后和 2024 年及之后。所有回归均包含职业、年份和行业固定效应，并加入“基线特征 $\times$ 年份”交互项，以满足条件平行趋势。回归采用岗位发布数量加权，标准误按 ISCO 职业码聚类。括号内为聚类稳健标准误。$^{***}p<0.01$，$^{**}p<0.05$，$^{*}p<0.10$。

\bigskip

\begin{longtable}{lcccc}
\caption{处理期安慰剂检验（仅冲击前样本）}\label{tab:app_placebo_timing} \\
\toprule
断点年份 & Shannon 熵 & 有效主题数 & 显著主题数 & HHI \\
\midrule
\endfirsthead
\toprule
断点年份 & Shannon 熵 & 有效主题数 & 显著主题数 & HHI \\
\midrule
\endhead
\bottomrule
\endfoot
2017 & $0.047^{**}$ (0.022) & $0.607$ (0.418) & $0.672^{**}$ (0.275) & $-0.068^{**}$ (0.029) \\
2018 & $0.037$ (0.026) & $0.277$ (0.503) & $0.580^{*}$ (0.321) & $-0.060^{*}$ (0.034) \\
2019 & $0.032$ (0.025) & $0.408$ (0.494) & $0.477$ (0.305) & $-0.047$ (0.031) \\
2020 & $0.042^{*}$ (0.023) & $0.862^{*}$ (0.500) & $0.536^{*}$ (0.277) & $-0.052^{*}$ (0.029) \\
2021 & $0.050^{**}$ (0.024) & $1.342^{**}$ (0.555) & $0.612^{**}$ (0.272) & $-0.054^{*}$ (0.030) \\
\midrule
2022（基准） & $0.054^{**}$ (0.022) & $1.191^{**}$ (0.506) & $0.652^{**}$ (0.253) & $-0.063^{**}$ (0.027) \\
\end{longtable}

\noindent 注：本表报告将处理期人为提前至冲击前各年份后的安慰剂检验结果，用于考察结论是否由任意时间断点驱动。若识别策略有效，则冲击前各伪断点不应产生稳定且显著的估计系数。所有回归均沿用主文基准设定，包括职业、年份和行业固定效应、条件平行趋势控制、岗位发布数量加权以及按 ISCO 职业码聚类的标准误。

\bigskip

\begin{longtable}{lll}
\caption{随机化推断与样条 DID 伪断点检验}\label{tab:app_ri_timing} \\
\toprule
检验 & 统计量 & 数值 \\
\midrule
\endfirsthead
\toprule
检验 & 统计量 & 数值 \\
\midrule
\endhead
\bottomrule
\endfoot
\multicolumn{3}{l}{\textit{A 组：随机化推断（500 次置换）}} \\[2pt]
基准估计 $\hat{\beta}$（Post $\geq$ 2022） & $\beta$ & $0.054$ \\
置换分布均值                                & $\bar{\beta}_{\text{perm}}$ & $0.0004$ \\
置换分布标准差                              & $\text{SD}_{\text{perm}}$   & $0.0032$ \\
随机化推断 $p$ 值（双侧）                  & $P(|\beta_{\text{perm}}| \geq |\hat{\beta}|)$ & $< 0.002$ \\
\midrule
\multicolumn{3}{l}{\textit{B 组：样条 DID 伪断点}} \\[2pt]
断点 = 2018 & $\beta_2$（加速项） & $0.001$~~($p = 0.958$) \\
断点 = 2019 & $\beta_2$（加速项） & $0.012$~~($p = 0.669$) \\
\end{longtable}

\noindent 注：本表汇总随机化推断和伪断点检验结果。随机化推断通过重复置换处理分配来评估基准估计值在零假设下的出现概率；伪断点检验则考察将结构性断点错误地设在 2018 或 2019 年时，是否仍会得到显著结果。若估计结果确实反映生成式 AI 冲击后的变化，则随机化推断的 $p$ 值应较低，而伪断点回归不应产生系统性显著结果。

\bigskip

\begin{figure}[htbp]
\centering
\includegraphics[width=0.9\textwidth]{../output/figures/event_study_entropy_cond_pt.pdf}
\caption{Shannon 熵的事件研究估计（基期 = 2021）。图中报告逐年交互项系数 $\hat{\beta}_t$ 及其 95\% 置信区间。空心圆为仅含三重固定效应的无控制规格，实心圆为加入基线特征 $\times$ 年份交互项的条件平行趋势规格。虚线标记处理期起点（2022 年），灰色阴影区域为冲击后时期。2022--2023 年系数接近零且不显著，显著效应集中在 2024 年，表明生成式 AI 对中国招聘市场技能需求结构的影响存在滞后传导。2025 年数据尚不完整，不宜做过多解读。}
\label{fig:app_event_study_timing}
\end{figure}

\bigskip

% ============================================================
\section{Rambachan \& Roth (2023) 敏感性分析}
% ============================================================

表~\ref{tab:app_honestdid} 报告了基于 Rambachan \& Roth (2023) Relative Magnitudes 框架的敏感性分析结果。该方法以冲击前各年系数的最大绝对偏离为参照基准，考察在允许冲击后趋势偏离达到参照基准的 $\bar{M}$ 倍时，目标参数的稳健置信区间是否仍排除零。分析基于条件平行趋势规格的事件研究系数及其聚类稳健方差—协方差矩阵，排除数据不完整的 2025 年。

\begin{table}[htbp]
\centering
\caption{Rambachan \& Roth (2023) 敏感性分析：Shannon 熵}\label{tab:app_honestdid}
\begin{threeparttable}
\begin{tabular}{lcccc}
\toprule
& \multicolumn{2}{c}{目标：2024 年效应} & \multicolumn{2}{c}{目标：2022--2024 等权平均} \\
\cmidrule(lr){2-3} \cmidrule(lr){4-5}
$\bar{M}$ & 稳健 CI & 含零？ & 稳健 CI & 含零？ \\
\midrule
0.00 & $[+0.001, +0.142]$ & 否  & $[-0.019, +0.069]$ & 是 \\
0.10 & $[-0.008, +0.168]$ & 是  & $[-0.027, +0.085]$ & 是 \\
0.25 & $[-0.040, +0.219]$ & 是  & $[-0.050, +0.120]$ & 是 \\
0.50 & $[-0.135, +0.332]$ & 是  & $[-0.114, +0.195]$ & 是 \\
1.00 & $[-0.345, +0.544]$ & 是  & $[-0.255, +0.336]$ & 是 \\
\bottomrule
\end{tabular}
\begin{tablenotes}[flushleft]
\small
\item 注：采用 Relative Magnitudes ($\Delta^{RM}$) 方法。$\bar{M}$ 为冲击后偏离相对于冲击前最大偏离的允许比例；$\bar{M} = 0$ 对应标准条件平行趋势假设。事前最大偏离为 $|\hat{\beta}_{2016}| = 0.075$。2024 年效应的击穿值为 $\bar{M} \approx 0.02$，即允许冲击后偏离达到冲击前最大偏离的约 2\% 时下界跨过零。等权平均在 $\bar{M} = 0$ 时已包含零，因为 2022--2023 年效应接近零。基于条件平行趋势规格的事件研究系数及 ISCO 聚类稳健方差—协方差矩阵，2025 年（数据不完整）未纳入分析。
\end{tablenotes}
\end{threeparttable}
\end{table}

\noindent 上述结果表明 DID 估计对条件平行趋势假设较为敏感。不过需要注意，$\bar{M}$ 的参照基准是最远端年份（2016 年）的残余偏离（$|\hat{\beta}_{2016}| = 0.075$），而 2017--2020 年的偏离已显著衰减（$|\hat{\beta}|$ 在 0.025--0.056 之间）。如果残余前趋势在时间上呈衰减态势，冲击后的实际偏离可能远小于以 2016 年为基准的保守参照，$\bar{M}$ 的有效值也相应更低。

\end{document}
